\sloppy
\chapter{Cơ sở lý thuyết}
\label{chap:c2}

Chương này trình bày toàn bộ nền tảng lý thuyết cốt lõi làm cơ sở cho việc nghiên cứu và phát triển phương pháp học động trong hệ thống kiểm tra lỗi cho màn hình điện thoại. Các nội dung được hệ thống hóa từ lý thuyết xử lý ảnh số, biểu diễn entropy không gian, kiến trúc mạng nơ-ron tích chập tự mã hóa, mô hình học dư ResNet, các thuật toán phân cụm không giám sát họ Neural Gas, đến lý thuyết phát hiện bất thường dựa trên phân phối thống kê và bài toán học liên tục trong môi trường công nghiệp.

\section{Tổng quan về Xử lý ảnh và Trích chọn đặc trưng}

Xử lý ảnh và trích chọn đặc trưng đóng vai trò tiền đề trong hệ thống thị giác máy tính phục vụ kiểm tra chất lượng sản phẩm công nghiệp. Trong bài toán kiểm định màn hình điện thoại, dữ liệu hình ảnh thu nhận thường chứa đựng sự biến đổi phức tạp về độ sáng, kết cấu bề mặt và các dạng khuyết tật có kích thước siêu nhỏ. Việc xây dựng một không gian biểu diễn dữ liệu tối ưu giúp làm nổi bật các tín hiệu bất thường, tạo điều kiện thuận lợi cho các mô hình học máy ở giai đoạn sau thực hiện phân loại và phân cụm chính xác. Mục này tập trung hệ thống hóa các nguyên lý cơ bản về biểu diễn ảnh số, lý thuyết thông tin Shannon và phương pháp tính toán entropy cục bộ trong phân tích kết cấu bề mặt màn hình.

\subsection{Biểu diễn ảnh số và Không gian mức xám}

Biểu diễn ảnh số là quá trình rời rạc hóa tín hiệu quang học liên tục thành ma trận các giá trị số trong không gian hai chiều \cite{gonzalez2018digital}. Trong không gian mức xám, mỗi điểm ảnh được đặc trưng bởi một giá trị cường độ sáng nằm trong khoảng xác định, thông thường là từ giá trị không đến hai trăm năm mươi lăm đối với định dạng tám bit. Cấu trúc ma trận ảnh $I$ có kích thước chiều cao $H$ và chiều rộng $W$ được mô tả toán học thông qua Công thức \ref{eq:image_matrix}:

\begin{equation}
I(x,y) \in [0, L-1], \quad \forall x \in \{1, 2, \dots, H\}, \quad y \in \{1, 2, \dots, W\}
\label{eq:image_matrix}
\end{equation}

Trong Công thức \ref{eq:image_matrix}, ký hiệu $I(x,y)$ đại diện cho cường độ sáng của điểm ảnh tại tọa độ hàng $x$ và cột $y$, $H$ là số lượng hàng của ma trận ảnh, $W$ là số lượng cột của ma trận ảnh, và $L$ đại diện cho số lượng mức xám tối đa của hệ thống rời rạc hóa.

Đối với hệ thống kiểm tra màn hình điện thoại, việc chỉ sử dụng dữ liệu mức xám hoặc dữ liệu màu ba kênh truyền thống thường gặp hạn chế khi xuất hiện các lỗi có độ tương phản thấp. Các khuyết tật như đốm tròn mờ, hạt bụi nhỏ hoặc vùng phân bố độ sáng không đều thường hòa lẫn vào kết cấu nền của màn hình. Để giải quyết thách thức này, phương pháp đề xuất mở rộng cấu trúc biểu diễn dữ liệu bằng cách bổ sung một kênh thông tin thứ tư chứa đặc trưng entropy cục bộ, tạo thành một tensor dữ liệu bốn kênh với cấu trúc $I_{\text{ext}} \in \mathbb{R}^{H \times W \times 4}$. Việc kết hợp ba kênh màu cơ bản hoặc kênh mức xám cùng kênh entropy giúp hệ thống duy trì đầy đủ thông tin hình thái hình học đồng thời tăng cường độ nhạy đối với các biến đổi kết cấu bất thường.

Cấu trúc biểu diễn của tensor dữ liệu mở rộng được mô tả trực quan thông qua Hình \ref{fig:four_channel_structure}.

\begin{figure}[h!]
\centering
\begin{tikzpicture}[every node/.style={transform shape, font=\normalsize}]
    % Channel 1
    \draw[fill=red!10, draw=red!60, thick] (0,0) -- (4,0) -- (4,3) -- (0,3) -- cycle;
    \node at (2,0.5) {Kênh Red};
    
    % Channel 2
    \draw[fill=green!15, draw=green!60, thick] (1,1) -- (5,1) -- (5,4) -- (1,4) -- cycle;
    \node at (3,1.5) {Kênh Green};
    
    % Channel 3
    \draw[fill=blue!20, draw=blue!60, thick] (2,2) -- (6,2) -- (6,5) -- (2,5) -- cycle;
    \node at (4,2.5) {Kênh Blue};
    
    % Channel 4 (Entropy)
    \draw[fill=teal!25, draw=teal!80, ultra thick] (3,3) -- (7,3) -- (7,6) -- (3,6) -- cycle;
    \node[align=center] at (5,4.5) {Kênh Entropy\\cục bộ\\\textit{(Kênh bổ sung)}};
    
    % Annotations
    \draw[->, thick, black] (7, 4.5) -- (9, 4.5) node[right, font=\normalsize] {Tensor đầu vào $H \times W \times 4$};
\end{tikzpicture}
\caption{Cấu trúc tensor dữ liệu bốn kênh kết hợp kênh màu và kênh entropy cục bộ.}
\label{fig:four_channel_structure}
\end{figure}

Trong Hình \ref{fig:four_channel_structure}, ba lớp đệm phía dưới thể hiện các kênh biểu diễn hình ảnh thông thường, trong khi lớp đệm phía trên được tô màu khác biệt đại diện cho kênh entropy cục bộ được trích xuất từ bề mặt màn hình. Sự tích hợp này đảm bảo toàn bộ thông tin không gian và thông tin thống kê biến động được truyền tải đồng thời vào các mạng nơ-ron phía sau.

\subsection{Lý thuyết thông tin và Shannon Entropy}

Lý thuyết thông tin do Claude Shannon đề xuất cung cấp khuôn khổ toán học định lượng độ hỗn loạn và lượng thông tin chứa đựng trong một nguồn phát ngẫu nhiên \cite{shannon1948mathematical}. Trong lĩnh vực xử lý ảnh số, một bức ảnh hoặc một vùng ảnh có thể xem là một nguồn phát thông tin rời rạc, trong đó các giá trị mức xám của điểm ảnh đại diện cho các trạng thái có thể xảy ra của biến ngẫu nhiên. Công thức tính Shannon Entropy được xác định theo Công thức \ref{eq:shannon_entropy}:

\begin{equation}
H(X) = -\sum_{i=1}^{N} P(x_i) \log_2 P(x_i)
\label{eq:shannon_entropy}
\end{equation}

Trong Công thức \ref{eq:shannon_entropy}, ký hiệu $H(X)$ thể hiện giá trị entropy của biến ngẫu nhiên rời rạc $X$, $x_i$ đại diện cho trạng thái thứ $i$ trong tập hợp các trạng thái có thể xảy ra, $P(x_i)$ là xác suất xuất hiện của trạng thái $x_i$, và $N$ thể hiện tổng số lượng trạng thái phân biệt của biến ngẫu nhiên $X$. Cơ số hai của phép toán lôgarit quy định đơn vị đo lượng thông tin tính bằng bit.

Khi áp dụng vào ảnh số, xác suất $P(r_k)$ của một mức xám $r_k$ được ước lượng thông qua biểu đồ tần suất không gian của toàn bộ bức ảnh hoặc cửa sổ xét duyệt. Công thức xác định xác suất $P(r_k)$ được thể hiện trong Công thức \ref{eq:histogram_probability}:

\begin{equation}
P(r_k) = \frac{n_k}{M \cdot N_w}
\label{eq:histogram_probability}
\end{equation}

Trong Công thức \ref{eq:histogram_probability}, ký hiệu $n_k$ đại diện cho số lượng điểm ảnh có mức xám bằng $r_k$, $M$ thể hiện chiều cao tính bằng điểm ảnh của vùng ảnh xét duyệt, và $N_w$ đại diện cho chiều rộng tính bằng điểm ảnh của vùng ảnh xét duyệt. Tích số giữa $M$ và $N_w$ thể hiện tổng số điểm ảnh chứa trong vùng không gian đang phân tích.

Giá trị Shannon Entropy đạt giá trị cực đại khi tất cả các mức xám trong vùng xét duyệt có xác suất xuất hiện đồng đều, phản ánh một vùng bề mặt có mức độ hỗn loạn cao hoặc chứa cấu trúc dữ liệu phức tạp. Ngược lại, nếu vùng ảnh hoàn toàn đồng nhất về mức xám, xác suất của mức xám đó bằng một và xác suất các mức xám khác bằng không, dẫn đến giá trị entropy bằng không. Tính chất này tạo tiền đề quan trọng cho việc phân biệt giữa vùng màn hình hiển thị bình thường và các vùng bị lỗi bề mặt.

\subsection{Entropy cục bộ trong phân tích kết cấu}

Mặc dù Shannon Entropy toàn cục cung cấp chỉ số đánh giá tổng quan về độ hỗn loạn của toàn bộ bức ảnh, chỉ số này không thể định vị được vị trí không gian của các biến đổi bất thường cục bộ \cite{haralick1973textural}. Phân tích kết cấu dựa trên entropy cục bộ giải quyết hạn chế này bằng cách áp dụng kỹ thuật cửa sổ trượt quét qua từng điểm ảnh trên ma trận hình ảnh. Tại mỗi vị trí điểm ảnh $(x,y)$, một vùng lân cận kích thước $w \times w$ được khởi tạo để tính toán giá trị phân bố xác suất xám nội vùng.

Công thức xác định entropy cục bộ tại tọa độ $(x,y)$ được biểu diễn chi tiết qua Công thức \ref{eq:local_entropy}:

\begin{equation}
H_{\text{local}}(x,y) = -\sum_{k=0}^{L-1} P_{W_{x,y}}(k) \log_2 \left( P_{W_{x,y}}(k) + \epsilon \right)
\label{eq:local_entropy}
\end{equation}

Trong Công thức \ref{eq:local_entropy}, ký hiệu $H_{\text{local}}(x,y)$ đại diện cho giá trị entropy cục bộ được gán cho điểm ảnh trung tâm tại tọa độ $(x,y)$, $W_{x,y}$ thể hiện ma trận cửa sổ trượt kích thước $w \times w$ nhận điểm ảnh $(x,y)$ làm tâm, $P_{W_{x,y}}(k)$ là xác suất xuất hiện của mức xám $k$ tính riêng trong phạm vi cửa sổ $W_{x,y}$, $L$ đại diện cho tổng số mức xám của hình ảnh, và $\epsilon$ là một hằng số dương siêu nhỏ được thêm vào nhằm tránh lỗi toán học khi tính lôgarit của giá trị không.

Quy trình trượt cửa sổ và chuyển đổi giá trị mức xám thành bản đồ entropy cục bộ được minh họa rõ ràng trong Hình \ref{fig:local_entropy_pipeline}.

\begin{figure}[h!]
\centering
\begin{tikzpicture}[every node/.style={transform shape, font=\normalsize}]
    % Input Image Grid (dịch sang trái 0.5)
    \draw[step=0.5cm, gray!40, thin] (-0.5,0) grid (3,3.5);
    \draw[black, thick] (-0.5,0) rectangle (3,3.5);
    \node[above, font=\normalsize] at (1.25,3.5) {Ma trận ảnh mức xám $I$};
    
    % Sliding Window
    \draw[fill=gray!30, draw=gray!80, thick] (0.5,1.5) rectangle (2,3);
    \node[font=\normalsize] at (1.25,2.25) {$W_{x,y}$};
    
    % Arrow to Calculation (dịch theo)
    \draw[->, thick, gray] (3.2, 1.75) -- (6.3, 1.75);
    \node[above, font=\normalsize] at (4.75,1.75) {Tính $P_{W}(k)$};
    
    % Output Entropy Map Grid (dịch sang phải thêm 1 đơn vị)
    \draw[step=0.5cm, gray!40, thin] (6.5,0) grid (10,3.5);
    \draw[black, thick] (6.5,0) rectangle (10,3.5);
    \node[above, font=\normalsize] at (8.25,3.5) {Bản đồ Entropy cục bộ};
    
    % Output Pixel Highlight
    \draw[fill=teal!40, draw=teal!90, thick] (8,2) rectangle (8.5,2.5);
\end{tikzpicture}
\caption{Quy trình tính toán bản đồ entropy cục bộ thông qua cửa sổ trượt lân cận.}
\label{fig:local_entropy_pipeline}
\end{figure}



Trong Hình \ref{fig:local_entropy_pipeline}, cửa sổ trượt $W_{x,y}$ di chuyển lần lượt qua từng vị trí trên ma trận ảnh gốc $I$. Giá trị entropy tính toán từ cửa sổ này được gán trực tiếp vào vị trí tương ứng trên bản đồ entropy cục bộ, tạo ra một ma trận hai chiều có cùng kích thước không gian với ảnh gốc.

Đối với bài toán kiểm tra lỗi màn hình điện thoại, sự hiện diện của các đốm tròn, vệt loang, nứt vỡ hoặc điểm chết màn hình sẽ làm đứt gãy tính đồng nhất của mô hình kết cấu xung quanh. Việc áp dụng tính toán entropy cục bộ mang lại ba ưu thế kĩ thuật quan trọng:
(i) tăng cường độ nổi bật của các đường biên và các chi tiết bất thường bất kể mức xám tuyệt đối của vùng lỗi;
(ii) giảm thiểu sự phụ thuộc vào điều kiện chiếu sáng không đều trên bề mặt hiển thị của sản phẩm;
(iii) chuẩn hóa thông tin nhiễu nền, giúp các mạng nơ-ron tự động nén đặc trưng thu nhận được tín hiệu lỗi một cách tập trung.

Bản đồ entropy cục bộ thu được sau quá trình biến đổi này được kết hợp trực tiếp vào ma trận ảnh ban đầu để tạo thành kênh thông tin thứ tư. Kênh dữ liệu bổ sung này giữ vai trò là thành phần định hướng không gian, giúp mô hình Autoencoder trong Pha một nén dữ liệu chính xác và giúp mô hình ResNet trong Pha hai trích xuất các đặc trưng hình thái sắc nét hơn đối với các dạng lỗi bề mặt phức tạp.
\section{Mạng Nơ-ron Tích chập và Autoencoder}

Mạng nơ-ron tích chập và mô hình Autoencoder đóng vai trò nền tảng trong việc tự động hóa quá trình trích xuất đặc trưng và nén dữ liệu hình ảnh. Trong bài toán phát hiện bất thường trên bề mặt màn hình điện thoại ở Pha 1, việc xử lý trực tiếp tensor ảnh bốn kênh có chiều không gian lớn gây ra chi phí tính toán cao và nhiễu dữ liệu. Bằng cách kết hợp khả năng học đặc trưng không gian cục bộ của phép tích chập và khả năng biểu diễn nén của Autoencoder, hệ thống chuyển đổi dữ liệu ảnh đầu vào thành các vectơ biểu diễn ẩn có số chiều nhỏ hơn nhưng bảo toàn các thông tin cấu trúc cốt lõi.

\subsection{Mạng Nơ-ron Tích chập}

Mạng nơ-ron tích chập là một lớp mô hình học sâu được thiết kế chuyên biệt để xử lý dữ liệu cấu trúc lưới hai chiều như ảnh số \cite{lecanu2015deep}. Nguyên lý hoạt động cốt lõi của mạng dựa trên ba cơ chế toán học cơ bản bao gồm: (i) trường thụ cảm cục bộ giúp mạng tập trung vào các vùng điểm ảnh lân cận; (ii) chia sẻ trọng số giúp giảm số lượng tham số cần huấn luyện và tạo tính bất biến đối với phép dịch chuyển; (iii) lấy mẫu giảm nhằm thu gọn kích thước không gian của bản đồ đặc trưng.

Phép toán tích chập hai chiều giữa ma trận ảnh đầu vào hoặc bản đồ đặc trưng lớp trước với một hạt nhân tích chập được định nghĩa theo Công thức \ref{eq:convolution_operation}:

\begin{equation}
S(i,j) = (I * K)(i,j) = \sum_{m=-M}^{M} \sum_{n=-N}^{N} I(i-m, j-n) K(m,n)
\label{eq:convolution_operation}
\end{equation}

Trong Công thức \ref{eq:convolution_operation}, ký hiệu $S(i,j)$ đại diện cho giá trị tại tọa độ $(i,j)$ trên bản đồ đặc trưng đầu ra, $I$ thể hiện ma trận dữ liệu đầu vào, $K$ là ma trận hạt nhân tích chập kích thước $(2M+1) \times (2N+1)$, $m$ và $n$ lần lượt là chỉ số không gian của hạt nhân tích chập theo chiều cao và chiều rộng, toán tử $*$ thể hiện phép tích chập không gian hai chiều.

Sau phép tích chập, một hàm kích hoạt phi tuyến như hàm đơn vị tuyến tính sửa sai được áp dụng để tăng năng lực biểu diễn các mối quan hệ phức tạp. Tiếp theo, các lớp lấy mẫu giảm, phổ biến nhất là lấy mẫu cực đại, được sử dụng để trích xuất các đặc trưng nổi bật nhất trong từng vùng lân cận. Quy trình trích xuất đặc trưng cơ bản thông qua tích chập và lấy mẫu được mô tả trong Hình \ref{fig:cnn_pipeline}.

\begin{figure}[h!]
\centering
\begin{tikzpicture}[every node/.style={transform shape, font=\normalsize}]
    % Input Layer (dịch sang trái 1.0)
    \draw[fill=white, draw=black, thick] (-1.6,0) rectangle (1.7,2.5);
    \node[align=center] at (0.0,1.25) {Tensor đầu vào\\$H \times W \times 4$};
    
    % Arrow 1
    \draw[->, thick, gray] (1.7,1.25) -- (3.7,1.25);
    \node[above, font=\normalsize] at (2.7,1.25) {Tích chập};
    
    % Feature Maps Layer (giữ nguyên, ở giữa)
    \draw[fill=white, draw=black, thick] (3.9,0.2) rectangle (5.9,2.2);
    \draw[fill=white, draw=black, thick] (4.1,0.4) rectangle (6.1,2.4);
    \node[align=center] at (5.1,1.4) {Bản đồ\\đặc trưng};
    
    % Arrow 2
    \draw[->, thick, gray] (6.3,1.25) -- (8.3,1.25);
    \node[above, font=\normalsize] at (7.3,1.25) {Lấy mẫu};
    
    % Pooled Feature Maps (dịch sang phải thêm 1.0)
    \draw[fill=white, draw=black, thick] (8.5,0.5) rectangle (10.0,2.0);
    \node[align=center] at (9.25,1.25) {Bản đồ\\nén};
\end{tikzpicture}
\caption{Quy trình xử lý dữ liệu qua lớp tích chập và lớp lấy mẫu giảm.}
\label{fig:cnn_pipeline}
\end{figure}

Trong Hình \ref{fig:cnn_pipeline}, dữ liệu đầu vào chứa bốn kênh thông tin được biến đổi qua các bộ lọc tích chập để tạo thành chuỗi bản đồ đặc trưng. Các bản đồ đặc trưng này tiếp tục trải qua quá trình lấy mẫu giảm nhằm thu nhỏ kích thước không gian nhưng vẫn giữ lại các tín hiệu hình thái quan trọng cho công đoạn mã hóa tiếp theo.

\subsection{Kiến trúc Autoencoder}

Mạng mã hóa tự động (Autoencoder) là một kiến trúc mạng nơ-ron học không giám sát được thiết kế nhằm học biểu diễn ẩn của dữ liệu bằng cách tái tạo lại chính dữ liệu đầu vào ở đầu ra \cite{goodfellow2016deep}. Cấu trúc tổng thể của Autoencoder gồm hai phần chính: (i) bộ mã hóa thực hiện ánh xạ dữ liệu đầu vào không gian cao sang không gian ẩn có số chiều nhỏ hơn; (ii) bộ giải mã thực hiện ánh xạ từ không gian ẩn trở lại không gian ban đầu để khôi phục ảnh tái tạo.

Bộ mã hóa biến đổi tensor đầu vào $x \in \mathbb{R}^{H \times W \times 4}$ thành vectơ không gian ẩn $z \in \mathbb{R}^d$ thông qua chuỗi các lớp tích chập được mô tả toán học trong Công thức \ref{eq:encoder_func}:

\begin{equation}
z = f_{\theta}(x) = \sigma(W_e x + b_e)
\label{eq:encoder_func}
\end{equation}

Trong Công thức \ref{eq:encoder_func}, ký hiệu $z$ đại diện cho vectơ không gian ẩn thu được, $f_{\theta}$ là hàm ánh xạ của bộ mã hóa phụ thuộc vào tập tham số $\theta$, $x$ là tensor ảnh đầu vào bốn kênh, $W_e$ thể hiện ma trận trọng số của bộ mã hóa, $b_e$ là vectơ độ lệch của bộ mã hóa, và $\sigma$ đại diện cho hàm kích hoạt phi tuyến.

Bộ giải mã nhận vectơ không gian ẩn $z$ và thực hiện quá trình tái tạo lại dữ liệu đầu vào thành $\hat{x} \in \mathbb{R}^{H \times W \times 4}$ thông qua các phép tích chập chuyển vị, được biểu diễn qua Công thức \ref{eq:decoder_func}:

\begin{equation}
\hat{x} = g_{\phi}(z) = \sigma'(W_d z + b_d)
\label{eq:decoder_func}
\end{equation}

Trong Công thức \ref{eq:decoder_func}, ký hiệu $\hat{x}$ đại diện cho tensor ảnh đầu ra tái tạo, $g_{\phi}$ là hàm ánh xạ của bộ giải mã phụ thuộc vào tập tham số $\phi$, $z$ là vectơ không gian ẩn, $W_d$ thể hiện ma trận trọng số của bộ giải mã, $b_d$ là vectơ độ lệch của bộ giải mã, và $\sigma'$ đại diện cho hàm kích hoạt ở lớp xuất dữ liệu.

Sơ đồ tổng quan về kiến trúc mạng Autoencoder xử lý dòng dữ liệu ảnh bốn kênh được thể hiện trong Hình \ref{fig:autoencoder_architecture}.
\begin{figure}[h!]
\centering
\begin{tikzpicture}[every node/.style={transform shape, font=\fontsize{13pt}{16pt}\selectfont}]
    % 1. Đầu vào x (Rộng 3.0: từ -1 đến 2)
    \draw[fill=white, draw=black, thick] (-1,0) rectangle (2,3);
    \node[align=center] at (0.5,1.5) {Đầu vào $x$\\$H \times W \times 4$};
    
    % Mũi tên 1 (Khoảng cách 0.5: từ 2 đến 2.5)
    \draw[->, thick, black] (2,1.5) -- (2.5,1.5);
    
    % 2. Bộ mã hóa (Rộng 2.5: từ 2.5 đến 5.0)
    \draw[fill=white, draw=black, thick] (2.5,0) -- (5,0.8) -- (5,2.2) -- (2.5,3) -- cycle;
    \node[align=center] at (3.75,1.5) {Bộ mã hóa\\$f_{\theta}$};
    
    % Mũi tên 2 (Khoảng cách 0.5: từ 5.0 đến 5.5)
    \draw[->, thick, black] (5,1.5) -- (5.5,1.5);
    
    % 3. Latent z (Rộng 1.8: từ 5.5 đến 7.3)
    \draw[fill=white, draw=black, thick] (5.5,0.8) rectangle (7.3,2.2);
    \node[align=center] at (6.4,1.5) {Latent\\$z$};
    
    % Mũi tên 3 (Khoảng cách 0.5: từ 7.3 đến 7.8)
    \draw[->, thick, black] (7.3,1.5) -- (7.8,1.5);
    
    % 4. Bộ giải mã (Rộng 2.5: từ 7.8 đến 10.3)
    \draw[fill=white, draw=black, thick] (7.8,0.8) -- (10.3,0) -- (10.3,3) -- (7.8,2.2) -- cycle;
    \node[align=center] at (9.05,1.5) {Bộ giải mã\\$g_{\phi}$};
    
    % Mũi tên 4 (Khoảng cách 0.5: từ 10.3 đến 10.8)
    \draw[->, thick, black] (10.3,1.5) -- (10.8,1.5);
    
    % 5. Tái tạo x_hat (Rộng 3.0: từ 10.8 đến 13.8)
    \draw[fill=white, draw=black, thick] (10.8,0) rectangle (13.8,3);
    \node[align=center] at (12.3,1.5) {Tái tạo $\hat{x}$\\$H \times W \times 4$};
\end{tikzpicture}
\caption{Kiến trúc tổng thể của mạng Autoencoder nén và tái tạo ảnh màn hình.}
\label{fig:autoencoder_architecture}
\end{figure}


Trong Hình \ref{fig:autoencoder_architecture}, hình thang thu nhỏ biểu thị quá trình nén dữ liệu của bộ mã hóa nhằm trích xuất vectơ đặc trưng $z$ tại vùng nút thắt. Hình thang mở rộng tiếp theo đại diện cho quá trình khôi phục lại cấu trúc tensor ảnh ban đầu của bộ giải mã.

Trong quy trình thực nghiệm của Pha 1, mạng Autoencoder chỉ được huấn luyện ngoại tuyến trên tập dữ liệu ảnh màn hình bình thường. Nhờ đó, không gian ẩn $z$ học được cách biểu diễn hiệu quả các đặc trưng phân bố của sản phẩm đạt tiêu chuẩn chất lượng. Vectơ đặc trưng ẩn $z$ thu được từ bộ mã hóa sau đó được trích xuất trực tiếp để làm dữ liệu đầu vào cho mô hình Growing Neural Gas ngoại tuyến, phục vụ bài toán phân định ngưỡng bất thường ở giai đoạn tiếp theo.

\subsection{Hàm Mất mát Tái tạo}

Hàm mất mát tái tạo đo lường mức độ sai lệch giữa dữ liệu ảnh gốc và dữ liệu ảnh thu được sau quá trình giải mã. Việc lựa chọn hàm mất mát phù hợp quyết định trực tiếp đến khả năng bảo toàn chi tiết cấu trúc nền và độ nhạy của mô hình đối với các đặc trưng cục bộ. Trong luận văn này, sai số bình phương trung bình được sử dụng làm hàm mục tiêu chính để tối ưu hóa tập tham số của mạng Autoencoder \cite{hinton2006reducing}.

Công thức toán học của hàm mất mát sai số bình phương trung bình áp dụng cho dữ liệu tensor bốn kênh được biểu diễn qua Công thức \ref{eq:mse_loss}:

\begin{equation}
\mathcal{L}_{\text{MSE}}(x, \hat{x}) = \frac{1}{C \cdot H \cdot W} \sum_{c=1}^{C} \sum_{i=1}^{H} \sum_{j=1}^{W} \left( x(c,i,j) - \hat{x}(c,i,j) \right)^2
\label{eq:mse_loss}
\end{equation}

Trong Công thức \ref{eq:mse_loss}, ký hiệu $\mathcal{L}_{\text{MSE}}(x, \hat{x})$ thể hiện tổn thất tái tạo tính theo trung bình bình phương sai số, $C$ đại diện cho số lượng kênh dữ liệu đầu vào với giá trị cố định bằng bốn, $H$ là chiều cao ảnh tính bằng điểm ảnh, $W$ là chiều rộng ảnh tính bằng điểm ảnh, $x(c,i,j)$ là giá trị độ sáng hoặc giá trị entropy gốc tại kênh $c$, hàng $i$, cột $j$, và $\hat{x}(c,i,j)$ là giá trị tương ứng được tái tạo bởi bộ giải mã.

Quá trình huấn luyện mô hình Autoencoder thực hiện tối ưu hóa tập tham số $\theta$ và $\phi$ thông qua giải thuật lan truyền ngược nhằm giảm thiểu giá trị tổn thất theo Công thức \ref{eq:optimization_objective}:

\begin{equation}
(\theta^*, \phi^*) = \arg\min_{\theta, \phi} \mathcal{L}_{\text{MSE}}\left(x, g_{\phi}(f_{\theta}(x))\right)
\label{eq:optimization_objective}
\end{equation}

Trong Công thức \ref{eq:optimization_objective}, ký hiệu $\theta^*$ và $\phi^*$ lần lượt là tập hợp các trọng số tối ưu của bộ mã hóa và bộ giải mã thu được sau quá trình huấn luyện, $\arg\min$ đại diện cho phép toán tìm tập tham số làm tối thiểu hóa giá trị của hàm mất mát tái tạo $\mathcal{L}_{\text{MSE}}$.

Đối với các mẫu ảnh không chứa khuyết tật, giá trị mất mát tái tạo $\mathcal{L}_{\text{MSE}}$ duy trì ở mức rất thấp do mô hình đã được tối ưu hóa trên phân bố dữ liệu này. Khi đưa một ảnh có chứa lỗi bề mặt vào mô hình, do mạng chưa từng được học các dạng cấu trúc lỗi bất thường, bộ giải mã sẽ không thể khôi phục chính xác các vùng khuyết tật đó. Sự gia tăng đột biến của sai số tái tạo tại các vùng bị lỗi chính là cơ sở quan trọng giúp gia tăng tính phân biệt của vectơ không gian ẩn $z$, hỗ trợ mô hình Growing Neural Gas thực hiện phát hiện bất thường chính xác hơn.
\section{Mạng nơ-ron tích chập ResNet cho trích xuất đặc trưng hình thái}

Trích xuất đặc trưng hình thái đóng vai trò quyết định đối với hiệu quả phân cụm các dạng khuyết tật bề mặt màn hình điện thoại trong Pha hai. Trong khi mô hình Autoencoder ở Pha một tập trung vào việc học phân phối ảnh bình thường để tái tạo dữ liệu, mô hình ở Pha hai cần trích xuất được các thông tin có tính phân biệt cao về hình học, hình thái và đường biên của vết lỗi. Kiến trúc mạng học dư ResNet được lựa chọn nhờ khả năng trích xuất đặc trưng sâu mà không gặp phải hiện tượng suy giảm tín hiệu gradient. Mục này phân tích cơ sở lý thuyết của học dư, phương pháp cải tiến lớp đầu vào cho luồng dữ liệu bốn kênh và vai trò của đặc trưng hình thái trong quy trình phân cụm trực tuyến của thuật toán EAGNG.

\subsection{Cơ chế Học dư và Kiến trúc Residual Block}

Khi tăng chiều sâu của các mạng nơ-ron tích chập truyền thống, mô hình thường gặp phải hiện tượng tiêu biến gradient hoặc bùng nổ gradient, dẫn đến độ chính xác huấn luyện bị suy giảm \cite{he2016deep}. Kiến trúc ResNet giải quyết triệt để thách thức này bằng cách giới thiệu cơ chế kết nối tắt hay còn gọi là đường truyền đồng nhất. Thay vì buộc các lớp tích chập phải học trực tiếp một ánh xạ mục tiêu đầy đủ, cơ chế này cấu trúc lại mạng để các lớp nơ-ron chỉ cần học phần dư chênh lệch giữa đầu vào và đầu ra.

Giả sử ánh xạ gốc cần tối ưu hóa là $\mathcal{H}(x)$, kiến trúc học dư chuyển đổi bài toán thành việc tối ưu hóa hàm phần dư $\mathcal{F}(x) = \mathcal{H}(x) - x$. Khi đó, ánh xạ mục tiêu ban đầu được khôi phục theo Công thức \ref{eq:residual_learning}:

\begin{equation}
y = \mathcal{F}(x, \{W_i\}) + x
\label{eq:residual_learning}
\end{equation}

Trong Công thức \ref{eq:residual_learning}, ký hiệu $y$ đại diện cho vectơ đầu ra của khối học dư, $x$ thể hiện vectơ đầu vào của khối, $W_i$ đại diện cho tập ma trận trọng số của các lớp tích chập trong khối, và $\mathcal{F}(x, \{W_i\})$ thể hiện ánh xạ phần dư cần học qua các lớp xử lý phi tuyến.

Khi kích thước của đầu vào $x$ và đầu ra $\mathcal{F}$ trùng khớp nhau, phép cộng trực tiếp theo từng phần tử được thực hiện. Trong trường hợp kích thước không gian hoặc số lượng kênh thay đổi, một ma trận biến đổi tuyến tính $W_s$ được bổ sung vào đường truyền tắt để chuẩn hóa kích thước, như được mô tả trong Công thức \ref{eq:residual_projection}:

\begin{equation}
y = \mathcal{F}(x, \{W_i\}) + W_s x
\label{eq:residual_projection}
\end{equation}

In Công thức \ref{eq:residual_projection}, ký hiệu $W_s$ đại diện cho ma trận ma trận chiếu tuyến tính được áp dụng để điều chỉnh số lượng kênh hoặc kích thước không gian của vectơ đầu vào $x$ sao cho đồng nhất với đầu ra của hàm phần dư.

Cấu trúc hình học và luồng dữ liệu bên trong một khối học dư cơ bản được mô tả chi tiết trong Hình \ref{fig:residual_block_architecture}.

\begin{figure}[h!]
\centering
\begin{tikzpicture}[every node/.style={transform shape, font=\fontsize{13pt}{16pt}\selectfont}, shift={(2,0)}]
    % Main path nodes (kích thước lớn hơn, khoảng cách đều nhau)
    \node (input)   at (0,6.0) [draw, fill=white, draw=black, thick, minimum width=3.2cm, minimum height=1cm] {Đầu vào $x$};
    \node (conv1)   at (0,4.5) [draw, fill=white, draw=black, thick, minimum width=4.0cm, minimum height=1cm] {Lớp Tích chập};
    \node (relu1)   at (0,3.0) [draw, fill=white, draw=black, thick, minimum width=4.0cm, minimum height=0.9cm] {Hàm ReLU};
    \node (conv2)   at (0,1.5) [draw, fill=white, draw=black, thick, minimum width=4.0cm, minimum height=1cm] {Lớp Tích chập};
    \node (sum)     at (0,0.0) [circle, fill=white, draw=black, ultra thick, inner sep=3pt] {$+$};
    \node (relu2)   at (0,-1.5) [draw, fill=white, draw=black, thick, minimum width=4.0cm, minimum height=0.9cm] {Hàm ReLU};
    \node (output)  at (0,-3.0) [draw, fill=white, draw=black, thick, minimum width=3.2cm, minimum height=1cm] {Đầu ra $y$};

    % Main path arrows
    \draw[->, thick, black] (input) -- (conv1);
    \draw[->, thick, black] (conv1) -- (relu1);
    \draw[->, thick, black] (relu1) -- (conv2);
    \draw[->, thick, black] (conv2) -- (sum);
    \draw[->, thick, black] (sum) -- (relu2);
    \draw[->, thick, black] (relu2) -- (output);

    % Shortcut path
    \draw[->, ultra thick, black] (0,5.25) -- (3.8,5.25) -- (3.8,0.0) node[midway, right, font=\fontsize{13pt}{16pt}\selectfont] {Đường truyền tắt $x$} -- (sum);
\end{tikzpicture}
\caption{Cấu trúc khối học dư cơ bản với đường truyền tắt bảo toàn thông tin.}
\label{fig:residual_block_architecture}
\end{figure}


Trong Hình \ref{fig:residual_block_architecture}, dòng dữ liệu gốc $x$ được chia làm hai nhánh. Nhánh chính trải qua các phép biến đổi tích chập phi tuyến để trích xuất tín hiệu biến đổi, trong khi nhánh đường truyền tắt bỏ qua các lớp biến đổi và cộng trực tiếp tín hiệu ban đầu vào đầu ra trước khi qua hàm kích hoạt cuối cùng.

\subsection{Trích xuất đặc trưng hình thái cho luồng dữ liệu bốn kênh}

Các mô hình ResNet tiền huấn luyện trên tập dữ liệu chuẩn thường được thiết kế để tiếp nhận đầu vào ảnh màu ba kênh. Trong hệ thống kiểm tra màn hình điện thoại, để khai thác tối đa thông tin từ kênh entropy cục bộ đã được ghép nối ở giai đoạn tiền xử lý, kiến trúc của lớp tích chập đầu tiên trong mạng ResNet cần được điều chỉnh cấu trúc trọng số \cite{szegedy2016rethinking}.

Ma trận hạt nhân tích chập của lớp đầu tiên được mở rộng chiều không gian kênh từ ba lên bốn. Phép toán tích chập đối với tensor đầu vào bốn kênh $I_{\text{ext}} \in \mathbb{R}^{H \times W \times 4}$ tại bản đồ đặc trưng thứ $m$ được xác định theo Công thức \ref{eq:resnet_4channel_conv}:

\begin{equation}
F_m(i,j) = \sum_{c=1}^{4} \left( I_{\text{ext}, c} * K_{m,c} \right)(i,j) + b_m
\label{eq:resnet_4channel_conv}
\end{equation}

Trong Công thức \ref{eq:resnet_4channel_conv}, ký hiệu $F_m(i,j)$ đại diện cho giá trị đặc trưng tại vị trí không gian $(i,j)$ trên bản đồ đặc trưng thứ $m$, $I_{\text{ext}, c}$ thể hiện kênh thứ $c$ của tensor ảnh đầu vào bốn kênh, $K_{m,c}$ là hạt nhân tích chập tương ứng với kênh $c$ của bộ lọc thứ $m$, toán tử $*$ thể hiện phép tích chập không gian hai chiều, và $b_m$ đại diện cho giá trị độ lệch của bộ lọc thứ $m$.

Việc tích hợp đồng thời bốn kênh dữ liệu mang lại các lợi thế kỹ thuật trong quá trình học đặc trưng hình thái bao gồm:
(i) ba kênh màu ban đầu cung cấp các thông tin về độ tương phản, mức xám và màu sắc tổng thể của vùng lỗi;
(ii) kênh entropy cục bộ bổ sung các tín hiệu về sự đứt gãy kết cấu và biên giới không gian của vết khuyết tật;
(iii) sự kết hợp qua các khối học dư cho phép mạng tự động tổng hợp các đặc trưng cấp cao biểu diễn chính xác hình thái hình học của các dạng lỗi phức tạp.

\subsection{Vai trò của ResNet trong quy trình phân cụm trực tuyến EAGNG}

Trong Pha hai của hệ thống, các hình ảnh chứa lỗi sau khi được phát hiện ở Pha một sẽ được chuyển sang công đoạn phân cụm trực tuyến. Mạng ResNet đóng vai trò là bộ trích xuất đặc trưng hình thái cố định, biến đổi mỗi ảnh khuyết tật thành một vectơ đặc trưng không gian $z_{\text{shape}} \in \mathbb{R}^d$. Nhờ khả năng biểu diễn hình học vượt trội của mạng tích chập sâu, các vết lỗi có cùng bản chất hình thái sẽ được ánh xạ thành các điểm dữ liệu nằm gần nhau trong không gian đặc trưng $d$ chiều.

Vectơ đặc trưng $z_{\text{shape}}$ sau đó được đưa trực tiếp vào mạng EAGNG để thực hiện học và phân cụm trực tuyến theo thời gian thực. Sự kết hợp này mang lại hai ý nghĩa quan trọng cho bài toán kiểm tra chất lượng sản phẩm:
(i) làm giảm số chiều của dữ liệu đầu vào cho mạng EAGNG so với việc sử dụng trực tiếp điểm ảnh thô, giúp tối ưu hóa tốc độ tính toán khoảng cách không gian;
(ii) loại bỏ các yếu tố nhiễu nền không quan trọng, tập trung vào các đặc trưng cấu trúc như tỷ lệ khung hình, độ dốc biên và diện tích phân bố của vết lỗi, tạo điều kiện thuận lợi cho thuật toán EAGNG xây dựng cấu trúc đồ thị phân cụm chính xác.
\section{Mạng nơ-ron không giám sát và Growing Neural Gas}

Học không giám sát đóng vai trò then chốt trong việc khám phá cấu trúc ẩn và phân cụm dữ liệu mà không cần nhãn tiên lượng. Trong bài toán phân tích khuyết tật màn hình điện thoại, các thuật toán học không giám sát dạng phân vùng dựa trên đồ thị cho phép biểu diễn phân phối dữ liệu phức tạp trong không gian đặc trưng nhiều chiều. Các mô hình dựa trên nguyên lý học cạnh tranh như Mạng tự tổ chức, Khí nơ-ron và Khí nơ-ron tăng trưởng phản ánh quá trình tự thích ứng của tập hợp nút đối với mật độ dữ liệu đầu vào. Mục này trình bày tiến trình phát triển lý thuyết từ các mô hình mạng không gian cố định đến mạng tăng trưởng linh hoạt, làm cơ sở khoa học cho việc đề xuất mô hình cải tiến ở các chương sau.

\subsection{Mạng tự tổ chức Kohonen}

Mạng tự tổ chức được Kohonen đề xuất nhằm ánh xạ không gian dữ liệu nhiều chiều về không gian lưới thấp chiều mà vẫn bảo toàn quan hệ lân cận giữa các mẫu dữ liệu \cite{kohonen1982self}. Cấu trúc mạng bao gồm một lưới các nút nơ-ron được sắp xếp cố định trên mặt phẳng hai chiều, trong đó mỗi nút $i$ gắn liền với một vectơ trọng số $w_i \in \mathbb{R}^d$. Khi tiếp nhận một vectơ dữ liệu đầu vào $x \in \mathbb{R}^d$, mạng xác định nút chiến thắng $c(x)$ có khoảng cách Euclid nhỏ nhất tới $x$ theo Công thức \ref{eq:som_winner}:

\begin{equation}
c(x) = \arg\min_i \|x - w_i\|
\label{eq:som_winner}
\end{equation}

Trong Công thức \ref{eq:som_winner}, ký hiệu $c(x)$ đại diện cho chỉ số của nút chiến thắng, $x$ thể hiện vectơ dữ liệu đầu vào trong không gian $d$ chiều, $w_i$ đại diện cho vectơ trọng số của nút nơ-ron thứ $i$, và $\| \cdot \|$ biểu diễn phép đo khoảng cách Euclid trong không gian thực.

Sau khi xác định nút chiến thắng, trọng số của nút chiến thắng và các nút trong vùng lân cận không gian lưới được cập nhật theo Công thức \ref{eq:som_update}:

\begin{equation}
w_i(t+1) = w_i(t) + h_{ci}(t) \left[ x(t) - w_i(t) \right]
\label{eq:som_update}
\end{equation}

Trong Công thức \ref{eq:som_update}, ký hiệu $w_i(t)$ đại diện cho vectơ trọng số của nút thứ $i$ tại bước thời gian $t$, $x(t)$ thể hiện mẫu dữ liệu đầu vào tại bước $t$, và $h_{ci}(t)$ đại diện cho hàm lân cận không gian giữa nút chiến thắng $c$ và nút $i$. Hàm lân cận $h_{ci}(t)$ thường được xác định bởi hàm Gauss suy giảm theo thời gian và khoảng cách lưới giữa hai nút, được tính theo Công thức \ref{eq:som_neighborhood}:

\begin{equation}
h_{ci}(t) = \alpha(t) \cdot \exp\left( -\frac{\|r_c - r_i\|^2}{2\sigma^2(t)} \right)
\label{eq:som_neighborhood}
\end{equation}

Trong Công thức \ref{eq:som_neighborhood}, ký hiệu $\alpha(t)$ đại diện cho tốc độ học tại bước $t$, $r_c$ và $r_i$ lần lượt thể hiện tọa độ vị trí cố định của nút $c$ và nút $i$ trên lưới hai chiều, và $\sigma(t)$ đại diện cho bán kính vùng lân cận suy giảm theo thời gian.

Mô hình Mạng tự tổ chức sở hữu ưu thế về khả năng trực quan hóa dữ liệu cấu trúc cao về dạng lưới hai chiều dễ quan sát. Tuy nhiên, kiến trúc này bộc lộ những hạn chế kỹ thuật đáng kể khi áp dụng vào dữ liệu khuyết tật thực tế:
(i) số lượng nút nơ-ron và cấu trúc lưới cố định buộc người dùng phải xác định trước, làm giảm tính dẻo khi phân phối dữ liệu thay đổi;
(ii) sự ràng buộc không gian lưới cứng nhắc dễ dẫn đến tình trạng các nút bị kéo dãn vào các vùng không có dữ liệu;
(iii) tốc độ học và bán kính lân cận suy giảm dần về không khiến mạng mất khả năng học các mẫu dữ liệu mới xuất hiện ở giai đoạn sau.

\subsection{Mạng Khí nơ-ron}

Mạng Khí nơ-ron do Martinetz và Schulten đề xuất nhằm loại bỏ sự ràng buộc của cấu trúc lưới không gian cố định trong Mạng tự tổ chức \cite{martinetz1991neural}. Thay vì dựa vào vị trí hình học trên lưới, Mạng Khí nơ-ron sắp xếp thứ tự các nút nơ-ron dựa trên khoảng cách không gian đặc trưng thực tế tới mẫu dữ liệu đầu vào. Đối với mỗi mẫu $x$, mạng tính toán khoảng cách tới toàn bộ các nút và gán chỉ số xếp hạng $k_i(x, W) \in \{0, 1, \dots, N-1\}$, trong đó nút gần nhất có chỉ số bằng không.

Quy tắc cập nhật trọng số cho tất cả các nút nơ-ron trong Mạng Khí nơ-ron được thực hiện theo Công thức \ref{eq:ng_update}:

\begin{equation}
w_i(t+1) = w_i(t) + \epsilon(t) \cdot \exp\left( -\frac{k_i(x, W)}{\lambda(t)} \right) \left[ x(t) - w_i(t) \right]
\label{eq:ng_update}
\end{equation}

Trong Công thức \ref{eq:ng_update}, ký hiệu $w_i(t)$ đại diện cho vectơ trọng số của nút thứ $i$ tại bước $t$, $x(t)$ là mẫu dữ liệu đầu vào, $k_i(x, W)$ thể hiện thứ tự xếp hạng khoảng cách của nút $i$ so với mẫu $x$, $\epsilon(t)$ đại diện cho tốc độ học tại bước $t$, và $\lambda(t)$ thể hiện hằng số suy giảm phạm vi ảnh hưởng của thứ tự xếp hạng.

Mô hình Khí nơ-ron đạt được tốc độ hội tụ nhanh hơn và giảm thiểu sai số biểu diễn đáng kể so với Mạng tự tổ chức nhờ không bị giới hạn bởi lưới hai chiều cố định. Mặc dù vậy, thuật toán vẫn duy trì các điểm hạn chế cốt lõi bao gồm:
(i) tổng số lượng nút nơ-ron $N$ vẫn phải được thiết lập cố định từ trước;
(ii) chi phí tính toán tăng nhanh do phải sắp xếp thứ tự khoảng cách cho toàn bộ các nút tại mỗi bước học;
(iii) sự suy giảm theo thời gian của các tham số $\epsilon(t)$ và $\lambda(t)$ cản trở mạng thích ứng với luồng dữ liệu học liên tục.

\subsection{Mạng Khí nơ-ron tăng trưởng}

Mạng khí nơ-ron tăng trưởng do Fritzke đề xuất là bước phát triển quan trọng giúp giải quyết triệt để hạn chế về số lượng nút cố định của các mô hình tiền nhiệm \cite{fritzke1995growing}. Thuật toán bắt đầu với một tập hợp gồm hai nút ban đầu và tự động bổ sung nút mới hoặc xóa bỏ các kết nối không hiệu quả dựa trên việc tích lũy sai số cục bộ. Mạng duy trì một tập hợp các nút $A$ và một tập hợp các cạnh nối $C$ giữa các cặp nút để phản ánh cấu trúc đồng dạng hình học của dữ liệu.

Quy trình cập nhật của thuật toán Mạng Khí nơ-ron tăng trưởng khi tiếp nhận một mẫu dữ liệu $x$ được thực hiện theo chuỗi các bước toán học:

Đầu tiên, mạng xác định hai nút gần $x$ nhất, ký hiệu là nút chiến thắng thứ nhất $s_1$ và nút chiến thắng thứ hai $s_2$, thỏa mãn điều kiện khoảng cách theo Công thức \ref{eq:gng_winners}:

\begin{equation}
s_1 = \arg\min_{i \in A} \|x - w_i\|, \quad s_2 = \arg\min_{i \in A \setminus \{s_1\}} \|x - w_i\|
\label{eq:gng_winners}
\end{equation}

Trong Công thức \ref{eq:gng_winners}, ký hiệu $s_1$ và $s_2$ lần lượt đại diện cho chỉ số của nút chiến thắng thứ nhất và thứ hai trong tập hợp nút $A$, $x$ thể hiện mẫu dữ liệu đầu vào, và $w_i$ đại diện cho vectơ trọng số của nút $i$.

Thứ hai, mạng tăng tuổi của tất cả các cạnh nối xuất phát từ $s_1$. Nếu giữa $s_1$ và $s_2$ chưa có cạnh nối, mạng tạo một cạnh mới nối hai nút này và đặt tuổi bằng không. Sai số tích lũy của nút chiến thắng $s_1$ được cập nhật theo Công thức \ref{eq:gng_error_update}:

\begin{equation}
E_{s_1} \leftarrow E_{s_1} + \|x - w_{s_1}\|^2
\label{eq:gng_error_update}
\end{equation}

Trong Công thức \ref{eq:gng_error_update}, ký hiệu $E_{s_1}$ đại diện cho giá trị sai số tích lũy của nút $s_1$, $x$ thể hiện mẫu dữ liệu đầu vào, và $w_{s_1}$ đại diện cho vectơ trọng số của nút $s_1$.

Thứ ba, vị trí trọng số của nút chiến thắng $s_1$ và các nút lân cận $n$ có kết nối trực tiếp với $s_1$ được điều chỉnh dịch chuyển về phía mẫu $x$ theo Công thức \ref{eq:gng_weight_update_s1} và Công thức \ref{eq:gng_weight_update_neighbors}:

\begin{equation}
w_{s_1} \leftarrow w_{s_1} + \epsilon_b (x - w_{s_1})
\label{eq:gng_weight_update_s1}
\end{equation}

\begin{equation}
w_n \leftarrow w_n + \epsilon_n (x - w_n), \quad \forall n \in N_{s_1}
\label{eq:gng_weight_update_neighbors}
\end{equation}

Trong Công thức \ref{eq:gng_weight_update_s1} và Công thức \ref{eq:gng_weight_update_neighbors}, ký hiệu $\epsilon_b$ đại diện cho tốc độ học của nút chiến thắng, $\epsilon_n$ thể hiện tốc độ học của các nút lân cận, và $N_{s_1}$ đại diện cho tập hợp các nút có cạnh nối trực tiếp với nút $s_1$.

Thứ tư, sau một chu kỳ cố định $\lambda$ bước học, mạng phát hiện nút $q$ có sai số tích lũy lớn nhất và nút $f$ có sai số lớn nhất trong số các nút lân cận của $q$. Một nút mới $r$ được chèn vào giữa $q$ và $f$ với trọng số được khởi tạo theo Công thức \ref{eq:gng_new_node}:

\begin{equation}
w_r = \frac{1}{2}(w_q + w_f)
\label{eq:gng_new_node}
\end{equation}

Trong Công thức \ref{eq:gng_new_node}, ký hiệu $w_r$ đại diện cho vectơ trọng số của nút mới bổ sung, $w_q$ và $w_f$ lần lượt thể hiện vectơ trọng số của hai nút có sai số lớn nhất. Sau khi chèn nút $r$, cạnh nối giữa $q$ và $f$ bị xóa bỏ, đồng thời hai cạnh mới nối $r$ với $q$ và $r$ với $f$ được tạo lập.

Quá trình thêm nút mới và tái cấu trúc liên kết đồ thị trong Mạng Khí nơ-ron tăng trưởng được mô tả trực quan trong Hình \ref{fig:gng_insertion_mechanism}.

\begin{figure}[h!]
\centering
\begin{tikzpicture}[every node/.style={transform shape, font=\fontsize{13pt}{16pt}\selectfont}]
    % Subfigure A: Before insertion
    \begin{scope}[shift={(0,0)}]
        \node[draw, circle, fill=blue!20, thick, minimum size=0.8cm] (q) at (0,2) {$q$};
        \node[draw, circle, fill=blue!20, thick, minimum size=0.8cm] (f) at (3,2) {$f$};
        \node[draw, circle, fill=gray!20, thick, minimum size=0.8cm] (n1) at (-1,0.5) {};
        \node[draw, circle, fill=gray!20, thick, minimum size=0.8cm] (n2) at (4,0.5) {};
        
        \draw[thick, black] (q) -- (f) node[midway, above, font=\fontsize{13pt}{16pt}\selectfont] {Cạnh cũ};
        \draw[thick, gray] (q) -- (n1);
        \draw[thick, gray] (f) -- (n2);
        
        \node[align=center, font=\fontsize{13pt}{16pt}\selectfont] at (1.5,-0.8) {Cấu trúc mạng trước khi thêm nút};
    \end{scope}

    % Transition Arrow
    \draw[->, ultra thick, black] (4.3, 1.2) -- (5.5, 1.2) node[midway, above, font=\fontsize{13pt}{16pt}\selectfont] {Thêm nút $r$};

    % Subfigure B: After insertion
    \begin{scope}[shift={(7,0)}]
        \node[draw, circle, fill=blue!20, thick, minimum size=0.8cm] (q2) at (0,2) {$q$};
        \node[draw, circle, fill=teal!30, draw=teal!80, ultra thick, minimum size=0.8cm] (r) at (1.5,2) {$r$};
        \node[draw, circle, fill=blue!20, thick, minimum size=0.8cm] (f2) at (3,2) {$f$};
        \node[draw, circle, fill=gray!20, thick, minimum size=0.8cm] (n12) at (-1,0.5) {};
        \node[draw, circle, fill=gray!20, thick, minimum size=0.8cm] (n22) at (4,0.5) {};
        
        \draw[thick, teal!80] (q2) -- (r);
        \draw[thick, teal!80] (r) -- (f2);
        \draw[thick, gray] (q2) -- (n12);
        \draw[thick, gray] (f2) -- (n22);
        
        \node[align=center, font=\normalsize] at (1.5,-0.8) {Cấu trúc mạng sau khi chèn nút $r$};
    \end{scope}
\end{tikzpicture}
\caption{Cơ chế chèn nút mới và phân tách kết nối trong Mạng Khí nơ-ron tăng trưởng.}
\label{fig:gng_insertion_mechanism}
\end{figure}

Trong Hình \ref{fig:gng_insertion_mechanism}, nút $q$ và nút $f$ đại diện cho hai nút có sai số tích lũy lớn nhất. Khi thỏa mãn điều kiện chu kỳ $\lambda$, nút mới $r$ được chèn vào trung điểm không gian giữa $q$ và $f$, đồng thời kết nối cũ giữa $q$ và $f$ được thay thế bằng hai kết nối mới tới $r$.

Mặc dù Mạng Khí nơ-ron tăng trưởng sở hữu năng lực tự thích ứng cấu trúc đồ thị mạnh mẽ, mô hình này vẫn bộc lộ các điểm hạn chế nghiêm trọng khi triển khai trong hệ thống học trực tuyến kiểm tra vết lỗi màn hình:
(i) tần suất thêm nút mới phụ thuộc hoàn toàn vào ngưỡng chu kỳ cố định $\lambda$ mà không xét đến mật độ dữ liệu thực tế tại vùng không gian xét duyệt, dẫn đến nguy cơ bùng nổ số lượng nút tại các vùng dữ liệu dày đặc;
(ii) mạng thiếu cơ chế phân biệt giữa dữ liệu nhiễu ngẫu nhiên và vết lỗi thực sự, khiến các mẫu nhiễu đứng đơn lẻ cũng tích lũy sai số và kích hoạt quá trình sinh nút không cần thiết;
(iii) hiện tượng quên thảm họa xảy ra khi luồng dữ liệu lỗi xuất hiện theo chuỗi thời gian, do các nút cũ liên tục bị điều chỉnh trọng số về phía dữ liệu mới mà không có cơ chế đóng băng hoặc bảo tồn bằng vùng đệm bằng chứng;
(iv) việc kết nối các nút chỉ dựa trên khoảng cách giữa hai nút chiến thắng $s_1$ và $s_2$ dễ gây ra hiện tượng nối nhầm cạnh giữa hai cụm khuyết tật nằm gần nhau về mặt không gian nhưng khác biệt về thuộc tính hình thái.

Những hạn chế nêu trên tạo nền tảng lý thuyết trực tiếp cho việc phát triển mô hình EAGNG trong luận văn này, nơi các cơ chế tích lũy bằng chứng, kiểm soát mật độ và kiểm tra chồng lấp hình học bao được tích hợp nhằm nâng cao tính ổn định và độ chính xác phân cụm.
\section{Cơ sở toán học cho các cải tiến của thuật toán EAGNG}

Mạng Khí nơ-ron tăng trưởng truyền thống bộc lộ những hạn chế kỹ thuật khi vận hành trong môi trường học trực tuyến, bao gồm hiện tượng bùng nổ số lượng nút tại các vùng dữ liệu có mật độ cao và nguy cơ hợp nhất sai các cụm khuyết tật do nhiễu không gian. Thuật toán EAGNG được phát triển nhằm khắc phục các tồn tại này thông qua bốn cải tiến cốt lõi dựa trên nền tảng toán học vững chắc. Mục này phân tích các nguyên lý toán học làm cơ sở cho việc điều khiển tốc độ học theo tuổi nút, cơ chế hoãn chèn nút dựa trên bằng chứng tích lũy, đánh giá mật độ không gian cục bộ và kiểm soát liên kết đồ thị bằng không gian bao hình học.

\subsection{Lý thuyết suy giảm tốc độ học thích ứng theo vòng đời nút}

Sự cân bằng giữa tính dẻo và tính ổn định là bài toán trung tâm trong các mô hình học liên tục và học trực tuyến \cite{grossberg1987competitive}. Trong mạng Khí nơ-ron tăng trưởng gốc, tốc độ học và mức độ tích lũy sai số của các nút nơ-ron được duy trì cố định đối với nút chiến thắng bất chấp thời gian tồn tại của nút đó trong mạng. Điều này khiến cho các nút nơ-ron đã hội tụ vững chắc vẫn bị nhiễu loạn khi xuất hiện các mẫu dữ liệu ngoại lai đơn lẻ.

Thuật toán EAGNG bổ sung biến số tuổi nút $a_i \in \mathbb{N}$ cho mỗi nút $i$ trong tập hợp nút $A$. Hàm suy giảm tốc độ học thích ứng theo thời gian sống của nút được định nghĩa theo Công thức \ref{eq:age_learning_rate}:

\begin{equation}
\eta_i(a_i) = \frac{\eta_0}{1 + \alpha_a \cdot a_i}
\label{eq:age_learning_rate}
\end{equation}

Trong Công thức \ref{eq:age_learning_rate}, ký hiệu $\eta_i(a_i)$ đại diện cho tốc độ học thích ứng của nút $i$ tại độ tuổi $a_i$, $\eta_0$ thể hiện tốc độ học khởi tạo ban đầu, và $\alpha_a$ là hệ số suy giảm theo tuổi nút nhận giá trị dương.

Tương tự, mức độ tích lũy sai số $E_i$ khi tiếp nhận mẫu dữ liệu mới $x$ được điều chỉnh suy giảm theo độ trưởng thành của nút, được thể hiện trong Công thức \ref{eq:age_error_accumulation}:

\begin{equation}
\Delta E_i = \frac{\|x - w_i\|^2}{1 + \beta_a \cdot a_i}
\label{eq:age_error_accumulation}
\end{equation}

Trong Công thức \ref{eq:age_error_accumulation}, ký hiệu $\Delta E_i$ đại diện cho giá trị gia tăng sai số tích lũy của nút $i$, $x$ thể hiện vectơ đặc trưng đầu vào, $w_i$ đại diện cho vectơ trọng số của nút $i$, $\| \cdot \|$ là phép đo khoảng cách Euclid, và $\beta_a$ là hệ số kiểm soát mức độ suy giảm tích lũy sai số.

Quy luật suy giảm thích ứng này mang lại hai hiệu ứng toán học quan trọng:
(i) các nút nơ-ron mới khởi tạo có độ tuổi nhỏ sẽ sở hữu tốc độ học cao, cho phép nhanh chóng di chuyển về trung tâm của vùng dữ liệu khuyết tật mới xuất hiện;
(ii) các nút nơ-ron đã tồn tại qua nhiều bước học sẽ có độ tuổi lớn, dẫn đến tốc độ học và mức gia tăng sai số giảm dần, giúp bảo vệ cấu trúc tri thức đã học khỏi hiện tượng nhiễu loạn ngẫu nhiên.

\subsection{Cơ chế đệm bộ nhớ và Ước lượng mật độ không gian cục bộ}

Việc chèn nút mới ngay lập tức khi xuất hiện một mẫu dữ liệu nằm xa các nút hiện có khiến mạng Khí nơ-ron tăng trưởng rất nhạy cảm với nhiễu ngẫu nhiên \cite{silva2006boundary}. Thuật toán EAGNG giải quyết thách thức này bằng cách kết hợp vùng đệm bộ nhớ tích lũy bằng chứng cùng phương pháp ước lượng mật độ nhân không gian.

Khi một mẫu dữ liệu đầu vào $x$ có khoảng cách tới nút chiến thắng vượt quá ngưỡng bất thường nhưng chưa đủ điều kiện tạo cấu trúc liên kết ngay, mẫu dữ liệu này được đưa vào vùng đệm bộ nhớ tạm thời $\mathcal{B} = \{x_1, x_2, \dots, x_K\}$. Mật độ dữ liệu cục bộ $D(x)$ tại điểm $x$ trong vùng đệm được ước lượng thông qua hàm nhân Gauss theo Công thức \ref{eq:kernel_density_estimation}:

\begin{equation}
D(x) = \frac{1}{K \cdot h^d} \sum_{j=1}^{K} \exp\left( -\frac{\|x - x_j\|^2}{2h^2} \right)
\label{eq:kernel_density_estimation}
\end{equation}

Trong Công thức \ref{eq:kernel_density_estimation}, ký hiệu $D(x)$ đại diện cho giá trị ước lượng mật độ không gian cục bộ tại vị trí $x$, $K$ thể hiện số lượng mẫu dữ liệu hiện có trong vùng đệm bộ nhớ, $h$ là bán kính cửa sổ nhân Gauss, $d$ đại diện cho số chiều của không gian đặc trưng, và $x_j$ thể hiện mẫu dữ liệu thứ $j$ trong vùng đệm.

Sơ đồ cơ chế hoãn chèn nút dựa trên vùng đệm bằng chứng và huấn luyện cục bộ được minh họa chi tiết trong Hình \ref{fig:evidence_buffer_mechanism}.

\begin{figure}[h!]
\centering
\begin{tikzpicture}[
    node distance=1.0cm and 1.5cm,
    block/.style={rectangle, draw=black, fill=white, thick, minimum width=3.2cm, minimum height=1.4cm, align=center, font=\normalsize},
    decision/.style={diamond, draw=black, fill=white, thick, aspect=1.8, inner sep=1pt, align=center, font=\normalsize},
    buffer/.style={rectangle, draw=black, fill=white, thick, minimum width=3.2cm, minimum height=1.4cm, align=center, font=\normalsize},
    train/.style={rectangle, draw=black, fill=white, thick, minimum width=3.2cm, minimum height=1.4cm, align=center, font=\normalsize},
    myarrow/.style={->, >=Stealth, thick, black}
]

    % Định vị các khối theo chiều dọc
    \node[block] (input) {Luồng dữ liệu\\mới $x$};
    
    % Khối quyết định 1
    \node[decision, below=0.8cm of input] (dec1) {Mẫu dị biệt?\\$\|x - w_{s_1}\| > T_{p2}$\\hoặc $D(x) < D_{\text{th}}$};
    
    % Nhánh "Không" của dec1 (rẽ sang bên phải)
    \node[block, right=1.2cm of dec1] (normal) {Cập nhật trọng số\\nút $s_1$ và lân cận};
    
    % Nhánh "Có" của dec1 (đi xuống dưới)
    \node[buffer, below=1.2cm of dec1] (buf) {Vùng đệm $\mathcal{B}$\\\textit{(Lưu trữ mẫu $x$)}};
    
    % Khối quyết định 2
    \node[decision, below=1.0cm of buf] (dec2) {Đủ bằng chứng?\\$K \ge K_{\text{min}}$};
    
    % Nhánh "Có" của dec2 (đi tiếp xuống dưới)
    \node[train, below=1.2cm of dec2] (local) {Huấn luyện cục bộ\\Tạo nút \& cạnh mới};

    % Vẽ mũi tên và nhãn
    \draw[myarrow] (input) -- (dec1);
    
    % Nhánh dec1
    \draw[myarrow] (dec1) -- node[midway, above] {Không} (normal);
    \draw[myarrow] (dec1) -- node[midway, right] {Có} (buf);
    
    \draw[myarrow] (buf) -- (dec2);
    
    % Nhánh dec2
    \draw[myarrow] (dec2) -- node[midway, right] {Có} (local);
    
    % Nhánh Không của dec2 (vòng sang bên trái quay lại vùng đệm B)
    \draw[myarrow] (dec2.west) -- ++(-1.2,0) |- node[pos=0.25, left, align=center] {Không\\(Tiếp tục tích lũy)} (buf.west);

\end{tikzpicture}
\caption{Sơ đồ cơ chế tích lũy bằng chứng trong vùng đệm $\mathcal{B}$ và kích hoạt huấn luyện cục bộ.}
\label{fig:evidence_buffer_mechanism}
\end{figure}


Trong Hình \ref{fig:evidence_buffer_mechanism}, các mẫu dữ liệu mới chưa đạt ngưỡng mật độ $D_{\text{th}}$ sẽ không tạo nút ngay mà được chuyển vào vùng đệm bộ nhớ $\mathcal{B}$. Quá trình huấn luyện cục bộ để sinh cấu trúc nút và cạnh nối mới chỉ được kích hoạt khi số lượng bằng chứng tích lũy $K$ đạt vượt ngưỡng tối thiểu $K_{\text{min}}$.

Bên cạnh đó, để tránh hiện tượng bùng nổ số lượng nút tại các vùng đã có mật độ biểu diễn cao, thuật toán EAGNG bổ sung chỉ số kiểm tra mật độ mạng cục bộ $\rho_i$ tại nút $i$ theo Công thức \ref{eq:node_density_check}:

\begin{equation}
\rho_i = \frac{1}{|N_i|} \sum_{j \in N_i} \frac{1}{\|w_i - w_j\| + \epsilon}
\label{eq:node_density_check}
\end{equation}

Trong Công thức \ref{eq:node_density_check}, ký hiệu $\rho_i$ đại diện cho mật độ không gian cục bộ xung quanh nút $i$, $N_i$ thể hiện tập hợp các nút lân cận có kết nối trực tiếp với nút $i$, $|N_i|$ là số lượng nút lân cận, $w_i$ và $w_j$ lần lượt là vectơ trọng số của nút $i$ và nút $j$, và $\epsilon$ là hằng số dương siêu nhỏ để tránh lỗi chia cho giá trị không.

Khi nút $q$ có sai số tích lũy lớn nhất đạt điều kiện phân tách, hệ thống kiểm tra chỉ số $\rho_q$. Việc chèn nút mới chỉ diễn ra khi giá trị $\rho_q$ nhỏ hơn ngưỡng mật độ tối đa $\rho_{\text{max}}$, đảm bảo mạng duy trì mật độ nút phân bố đồng đều và tối ưu tài nguyên tính toán.

\subsection{Lý thuyết hình học bao không gian trong kiểm soát liên kết đồ thị}

Trong quá trình học trực tuyến, các vết khuyết tật khác nhau có thể xuất hiện gần nhau trong không gian đặc trưng. Mô hình Khí nơ-ron tăng trưởng truyền thống chỉ dựa vào khoảng cách giữa hai nút chiến thắng $s_1$ và $s_2$ để tạo cạnh nối, dẫn đến rủi ro nối nhầm các cụm khuyết tật riêng biệt thành một phân lớp duy nhất \cite{he2008graph}.

Thuật toán EAGNG tích hợp cơ chế kiểm soát liên kết dựa trên lý thuyết không gian bao hình học chữ nhật song song với các trục tọa độ. Đối với mỗi cụm nơ-ron $C_m$ trong đồ thị, hình học bao không gian $B(C_m)$ được xác định bởi tập hợp các khoảng giá trị cực trị trên từng chiều không gian theo Công thức \ref{eq:bounding_box_def}:

\begin{equation}
B(C_m) = \prod_{d=1}^{D} \left[ \min_{i \in C_m} w_{i,d}, \, \max_{i \in C_m} w_{i,d} \right]
\label{eq:bounding_box_def}
\end{equation}

Trong Công thức \ref{eq:bounding_box_def}, ký hiệu $B(C_m)$ đại diện cho hình học bao không gian của cụm thứ $m$, $D$ là tổng số chiều của không gian đặc trưng, $w_{i,d}$ thể hiện giá trị tọa độ chiều thứ $d$ của nút $i$ thuộc cụm $C_m$, và toán tử $\prod$ biểu thị tích Cartesian của các khoảng giá trị không gian.

Khi xuất hiện yêu cầu tạo cạnh nối giữa nút $u \in C_m$ và nút $v \in C_n$ thuộc hai cụm phân biệt, thuật toán tiến hành tính toán thể tích vùng giao cắt $V_{\text{overlap}}$ giữa hai hình học bao $B(C_m)$ và $B(C_n)$ theo Công thức \ref{eq:bounding_box_overlap}:

\begin{equation}
V_{\text{overlap}}(C_m, C_n) = \prod_{d=1}^{D} \max\left(0, \, \min\left(w_{m,d}^{\text{max}}, w_{n,d}^{\text{max}}\right) - \max\left(w_{m,d}^{\text{min}}, w_{n,d}^{\text{min}}\right)\right)
\label{eq:bounding_box_overlap}
\end{equation}

Trong Công thức \ref{eq:bounding_box_overlap}, ký hiệu $V_{\text{overlap}}(C_m, C_n)$ đại diện cho thể tích phần không gian chồng lấp giữa hai hình học bao, $w_{m,d}^{\text{max}}$ và $w_{m,d}^{\text{min}}$ lần lượt là giá trị tọa độ lớn nhất và nhỏ nhất theo chiều $d$ của cụm $C_m$.

Tỷ lệ giao cắt không gian $R_{\text{overlap}}$ so với thể tích nhỏ hơn giữa hai không gian bao được xác định theo Công thức \ref{eq:overlap_ratio}:

\begin{equation}
R_{\text{overlap}} = \frac{V_{\text{overlap}}(C_m, C_n)}{\min\left(V(B(C_m)), \, V(B(C_n))\right)}
\label{eq:overlap_ratio}
\end{equation}

Trong Công thức \ref{eq:overlap_ratio}, ký hiệu $R_{\text{overlap}}$ đại diện cho tỷ lệ chồng lấp không gian giữa hai cụm, và $V(B(C_m))$ thể hiện tổng thể tích không gian bao của cụm $C_m$.

Quy tắc quyết định liên kết đồ thị được quy định cụ thể: nếu tỷ lệ giao cắt $R_{\text{overlap}}$ vượt quá ngưỡng cho phép $\tau_{\text{overlap}}$, hai cụm được xác nhận là thuộc cùng một vùng khuyết tật mở rộng và việc tạo cạnh nối được phê duyệt; ngược lại, nếu tỷ lệ $R_{\text{overlap}}$ nhỏ hơn ngưỡng, phép tạo cạnh nối bị từ chối nhằm duy trì sự ranh giới độc lập giữa hai dạng lỗi bề mặt. Cơ chế này đảm bảo mạng EAGNG phân cụm chính xác hình thái các vết khuyết tật phức tạp mà không bị sáp nhập sai lệch.
\section{Phát hiện bất thường và phân cụm dựa trên lý thuyết đồ thị}

Hệ thống kiểm tra chất lượng màn hình điện thoại đòi hỏi sự phối hợp chặt chẽ giữa khả năng sàng lọc mẫu bất thường ở Pha một và phân loại khuyết tật ở Pha hai. Việc kết hợp nguyên lý phân phối thống kê cùng lý thuyết đồ thị tạo nên khuôn khổ toán học hoàn chỉnh cho phép phát hiện các khuyết tật vượt khỏi phân phối của ảnh bình thường, đồng thời gom nhóm tự động các dạng lỗi bề mặt có đặc trưng tương đồng. Mục này trình bày nguyên lý phát hiện bất thường dựa trên quy tắc ba Sigma và phương pháp phân cụm đồ thị bằng thành phần liên thông.

\subsection{Nguyên lý phát hiện bất thường dựa trên phân phối thống kê và khoảng cách 3-Sigma}

Phát hiện bất thường dựa trên phân phối thống kê dựa trên giả định rằng các mẫu dữ liệu bình thường tập trung trong một miền không gian ẩn xác định và tuân theo phân phối chuẩn. Trong Pha một của quy trình đề xuất, mạng Autoencoder thực hiện nén các tensor ảnh bốn kênh bình thường thành các vectơ không gian ẩn $z \in \mathbb{R}^d$, sau đó mạng Khí nơ-ron tăng trưởng ngoại tuyến học cấu trúc phân phối của tập hợp các vectơ này. Mỗi mẫu ảnh thử nghiệm $x_{\text{test}}$ khi đi qua bộ mã hóa sẽ thu được vectơ ẩn $z_{\text{test}}$, và khoảng cách đại diện cho sai số tái tổ chức tới nút nơ-ron chiến thắng $w_c$ trong mạng Khí nơ-ron tăng trưởng được xác định theo Công thức \ref{eq:distance_to_winner}:

\begin{equation}
d(z_{\text{test}}, w_c) = \| z_{\text{test}} - w_c \| = \sqrt{\sum_{k=1}^{d} \left( z_{\text{test}, k} - w_{c, k} \right)^2}
\label{eq:distance_to_winner}
\end{equation}

Trong Công thức \ref{eq:distance_to_winner}, ký hiệu $d(z_{\text{test}}, w_c)$ đại diện cho khoảng cách Euclid từ vectơ không gian ẩn thử nghiệm $z_{\text{test}}$ tới vectơ trọng số của nút chiến thắng $w_c$, $d$ là số chiều của không gian ẩn, $z_{\text{test}, k}$ là thành phần thứ $k$ của vectơ ẩn thử nghiệm, và $w_{c, k}$ đại diện cho thành phần thứ $k$ của nút chiến thắng $w_c$.

Giá trị trung bình và độ lệch chuẩn của khoảng cách được tính toán trên toàn bộ tập dữ liệu huấn luyện ảnh bình thường để thiết lập ngưỡng phân định bất thường khoa học \cite{hawkins1980identification}. Kỳ vọng toán học và độ lệch chuẩn của khoảng cách không gian ẩn được định nghĩa lần lượt trong Công thức \ref{eq:mean_distance} và Công thức \ref{eq:std_distance}:

\begin{equation}
\mu = \frac{1}{N_{\text{good}}} \sum_{i=1}^{N_{\text{good}}} d(z_i, w_{c(i)})
\label{eq:mean_distance}
\end{equation}

\begin{equation}
\sigma = \sqrt{\frac{1}{N_{\text{good}}} \sum_{i=1}^{N_{\text{good}}} \left( d(z_i, w_{c(i)}) - \mu \right)^2}
\label{eq:std_distance}
\end{equation}

Trong Công thức \ref{eq:mean_distance} và Công thức \ref{eq:std_distance}, ký hiệu $\mu$ đại diện cho giá trị khoảng cách trung bình của tập mẫu ảnh bình thường, $\sigma$ thể hiện độ lệch chuẩn của phân phối khoảng cách, $N_{\text{good}}$ đại diện cho tổng số lượng mẫu ảnh bình thường trong tập huấn luyện ngoại tuyến, $z_i$ là vectơ ẩn của mẫu thứ $i$, và $w_{c(i)}$ là nút chiến thắng tương ứng với mẫu $z_i$.

Quy tắc ba Sigma được áp dụng để xác định ranh giới phân định bất thường với xác suất tin cậy đạt mức chín mươi chín phẩy bảy mươi ba phần trăm đối với phân phối chuẩn \cite{pukelsheim1994three}. Điều kiện quyết định một mẫu dữ liệu ẩn có thuộc nhóm bất thường và cần chuyển sang Pha hai hay không được thể hiện trong Công thức \ref{eq:three_sigma_decision}:

\begin{equation}
\text{Anomaly}(z_{\text{test}}) = \begin{cases} 1, & \text{nếu } d(z_{\text{test}}, w_c) > \mu + 3\sigma \\ 0, & \text{ngược lại} \end{cases}
\label{eq:three_sigma_decision}
\end{equation}

Trong Công thức \ref{eq:three_sigma_decision}, ký hiệu $\text{Anomaly}(z_{\text{test}})$ nhận giá trị bằng một khi mẫu thử nghiệm bị kết luận là chứa khuyết tật bất thường, và nhận giá trị bằng không khi mẫu thử nghiệm nằm trong giới hạn biến động cho phép của sản phẩm đạt tiêu chuẩn.

Phân bố xác suất và ngưỡng phân định bất thường dựa trên khoảng cách ba Sigma được minh họa trực quan trong Hình \ref{fig:three_sigma_distribution}.

\begin{figure}[h!]
\centering
\begin{tikzpicture}[every node/.style={transform shape, font=\normalsize}]
    % 1. Tô màu Vùng bình thường (từ 0.2 đến mu + 3*sigma = 6.0)
    \fill[cyan!15] (0.2,0) -- plot[domain=0.2:6, samples=100] (\x, {3.5*exp(-(\x-3)*(\x-3)/2)}) -- (6,0) -- cycle;
    
    % 2. Tô màu Vùng bất thường (phần đuôi > mu + 3*sigma = 6.0)
    \fill[red!25] (6,0) -- plot[domain=6:7.5, samples=50] (\x, {3.5*exp(-(\x-3)*(\x-3)/2)}) -- (7.5,0) -- cycle;
    
    % 3. Đường cong Gauss (Đỉnh mu = 3, sigma = 1)
    \draw[thick, blue!80!black] plot[domain=0.2:7.5, samples=100] (\x, {3.5*exp(-(\x-3)*(\x-3)/2)});
    
    % 4. Trục tọa độ (Chỉ biểu diễn miền khoảng cách d >= 0)
    \draw[->, thick, black] (-1,0) -- (8.0,0) node[right, font=\normalsize] {Khoảng cách $d(z, w_c)$};
    \draw[->, thick, black] (-1,0) -- (-1,4.0) node[above, font=\normalsize] {Mật độ xác suất};
    
    % 5. Đường trung bình mu
    \draw[dashed, blue!90!black, thick] (3,0) -- (3,3.5);
    \node[below, font=\normalsize] at (3,-0.1) {$\mu$};
    
    % Các vạch chia sigma tham khảo
    \draw[gray] (1,0) -- (1,0.1) node[below=2pt, xshift=-0.3cm, font=\small] {$\mu-2\sigma$};
    \draw[gray] (2,0) -- (2,0.1) node[below=2pt, font=\small] {$\mu-\sigma$};
    \draw[gray] (4,0) -- (4,0.1) node[below=2pt, font=\small] {$\mu+\sigma$};
    \draw[gray] (5,0) -- (5,0.1) node[below=2pt, xshift=0.3cm, font=\small] {$\mu+2\sigma$};
    
    % 6. Vạch ngưỡng mu + 3*sigma
    \draw[ultra thick, red!80!black] (6,0) -- (6,2.6);
    \node[above, font=\normalsize, align=center, red!80!black] at (6,2.6) {Ngưỡng $\mu + 3\sigma$};
    
    % 7. Chú thích các vùng
    % Nhãn Vùng bình thường
    \node[align=center, font=\normalsize, fill=cyan!15, inner sep=3pt, rounded corners=2pt] at (3,1.5) {Vùng ảnh bình thường\\(Pha 1: Chấp nhận)};
    
    % Mũi tên chỉ vào Vùng bất thường ở phần đuôi
    \node[align=center, font=\normalsize, text=red!80!black] (anomaly_text) at (9, 1.4) {Vùng bất thường\\(Pha 2: Chuyển EAGNG)};
    \draw[->, >=stealth, thick, red!80!black] (anomaly_text.south) .. controls (6.8,0.7) .. (6.3, 0.12);

\end{tikzpicture}
\caption{Phân bố khoảng cách Gauss và ngưỡng phân định bất thường dựa trên quy tắc ba Sigma.}
\label{fig:three_sigma_distribution}
\end{figure}

Trong Hình \ref{fig:three_sigma_distribution}, đường cong biểu diễn mật độ xác suất khoảng cách của các mẫu ảnh đạt tiêu chuẩn. Vùng tô màu nhạt bên trái đại diện cho miền dữ liệu bình thường, trong khi ranh giới đánh dấu tại điểm ba Sigma tách biệt các mẫu bất thường cần được chuyển tiếp sang Pha hai để thực hiện phân cụm.

\subsection{Phân cụm đồ thị bằng thành phần liên thông}

Phân cụm dựa trên thành phần liên thông chuyển đổi cấu trúc đồ thị mạng nơ-ron thành các tập hợp con rời rạc đại diện cho từng loại khuyết tật độc lập. Sau khi mạng EAGNG hoàn tất quá trình huấn luyện trực tuyến ở Pha hai, cấu trúc mạng được biểu diễn dưới dạng một đồ thị vô hướng $G = (V, E)$, trong đó $V$ là tập hợp tất cả các nút nơ-ron đại diện cho các trung tâm dữ liệu khuyết tật, và $E$ là tập hợp các cạnh nối phản ánh mối quan hệ lân cận giữa các nút.

Hai nút nơ-ron $u \in V$ và $v \in V$ được coi là thuộc cùng một thành phần liên thông nếu tồn tại ít nhất một đường đi nối trực tiếp hoặc gián tiếp giữa hai nút đó trên đồ thị $G$ \cite{hopcroft1973algorithm}. Định nghĩa toán học của tập hợp thành phần liên thông thứ $k$, ký hiệu là $C_k$, được thể hiện qua Công thức \ref{eq:connected_component}:

\begin{equation}
C_k = \{ v \in V \mid \exists p = (u = v_0, v_1, \dots, v_m = v), \, \forall j \in \{0, \dots, m-1\}, \, (v_j, v_{j+1}) \in E \}
\label{eq:connected_component}
\end{equation}

Trong Công thức \ref{eq:connected_component}, ký hiệu $C_k$ đại diện cho tập hợp các nút thuộc thành phần liên thông thứ $k$, $p$ thể hiện dãy các nút tạo thành đường đi từ nút $u$ đến nút $v$, $m$ là độ dài tính bằng số cạnh của đường đi, và $(v_j, v_{j+1})$ đại diện cho cạnh nối trực tiếp giữa nút $v_j$ và nút $v_{j+1}$ nằm trong tập cạnh $E$.

Tập hợp tất cả các thành phần liên thông tạo thành một phép phân hoạch hoàn chỉnh trên tập nút $V$, thỏa mãn các điều kiện phủ kín và rời rạc theo Công thức \ref{eq:graph_partition}:

\begin{equation}
V = \bigcup_{k=1}^{M} C_k, \quad C_i \cap C_j = \emptyset, \quad \forall i \neq j
\label{eq:graph_partition}
\end{equation}

Trong Công thức \ref{eq:graph_partition}, ký hiệu $M$ đại diện cho tổng số lượng cụm khuyết tật độc lập được trích xuất từ đồ thị, $C_i$ và $C_j$ lần lượt thể hiện thành phần liên thông thứ $i$ và thứ $j$, và $\emptyset$ đại diện cho tập hợp rỗng.

Cấu trúc phân tách các cụm khuyết tật riêng biệt thông qua thuật toán tìm thành phần liên thông trên đồ thị EAGNG được minh họa trong Hình \ref{fig:connected_components_graph}.

\begin{figure}[h!]
\centering
\begin{tikzpicture}[every node/.style={transform shape, font=\normalsize}]
    % Cluster 1 (Scratch defects)
    \draw[dashed, teal!80, thick, fill=teal!5] (-1,1.5) ellipse (1.8cm and 1.2cm);
    \node[draw, circle, fill=blue!30, thick, inner sep=3pt] (n1) at (-1.7,1.8) {};
    \node[draw, circle, fill=blue!30, thick, inner sep=3pt] (n2) at (-0.8,2.1) {};
    \node[draw, circle, fill=blue!30, thick, inner sep=3pt] (n3) at (-0.7,1.2) {};
    \node[draw, circle, fill=blue!30, thick, inner sep=3pt] (n4) at (-1.4,1.0) {};
    
    \draw[thick, gray] (n1) -- (n2);
    \draw[thick, gray] (n2) -- (n3);
    \draw[thick, gray] (n3) -- (n4);
    \draw[thick, gray] (n4) -- (n1);
    \node[above, font=\normalsize] at (-1,2.8) {Cụm lỗi $C_1$ (Lỗi hình tròn)};

    % Cluster 2 (Stain defects)
    \draw[dashed, cyan!80, thick, fill=cyan!5] (6,1.5) ellipse (1.6cm and 1.1cm);
    \node[draw, circle, fill=teal!40, thick, inner sep=3pt] (m1) at (5.3,1.8) {};
    \node[draw, circle, fill=teal!40, thick, inner sep=3pt] (m2) at (6.3,1.9) {};
    \node[draw, circle, fill=teal!40, thick, inner sep=3pt] (m3) at (6.6,1.1) {};
    \node[draw, circle, fill=teal!40, thick, inner sep=3pt] (m4) at (5.7,1.0) {};
    
    \draw[thick, gray] (m1) -- (m2);
    \draw[thick, gray] (m2) -- (m3);
    \draw[thick, gray] (m3) -- (m4);
    \node[above, font=\normalsize] at (6,2.8) {Cụm lỗi $C_2$ (Lỗi hình chữ nhật)};

    % No connection between clusters annotation
    \draw[<->, thick, dashed, gray] (1,1.5) -- (4,1.5);
    \node[above, font=\normalsize] at (2.6,1.5) {Không có cạnh};
    \node[below, font=\normalsize] at (2.6,1.5) {($C_1 \cap C_2 = \emptyset$)};
\end{tikzpicture}
\caption{Phân cụm các thành phần liên thông độc lập đại diện cho các dạng khuyết tật bề mặt.}
\label{fig:connected_components_graph}
\end{figure}

Trong Hình \ref{fig:connected_components_graph}, các nút nơ-ron thuộc cùng một phân lớp khuyết tật duy trì các liên kết cạnh nội vùng, tạo thành các đường đi liên thông khép kín. Việc thiếu hụt cạnh nối giữa hai nhóm nút giúp phân tách hoàn toàn các dạng khuyết tật khác nhau thành các cụm độc lập.

Phương pháp phân cụm đồ thị bằng thành phần liên thông sở hữu ba ưu thế kỹ thuật vượt trội trong môi trường sản xuất thực tế:
(i) thuật toán tự động xác định số lượng cụm khuyết tật $M$ dựa trên cấu trúc đồ thị thực tế mà không đòi hỏi thiết lập trước số cụm như các thuật toán phân hoạch truyền thống;
(ii) khả năng biểu diễn các cụm có hình dạng phi tuyến và cấu trúc hình học phức tạp nhờ tính linh hoạt của các kết nối cạnh;
(iii) sự kết hợp giữa các thuật toán kiểm soát liên kết của EAGNG và thành phần liên thông giúp loại bỏ các cạnh nối không hợp lệ, ngăn ngừa hiện tượng gộp sai các dạng khuyết tật bề mặt có vị trí gần nhau.
\section{Bài toán Học liên tục}

Bài toán học liên tục đóng vai trò trung tâm trong các hệ thống học máy vận hành trên luồng dữ liệu biến đổi theo thời gian. Môi trường vận hành thực tế luôn tồn tại sự dịch chuyển phân bố dữ liệu đầu vào do các yếu tố vật lý và điều kiện môi trường. Các mô hình học máy truyền thống thường được thiết kế dựa trên giả định dữ liệu độc lập và phân bố đồng nhất, dẫn đến sự suy giảm hiệu năng khi áp dụng vào luồng dữ liệu mở. Lĩnh vực học liên tục tập trung nghiên cứu các nguyên lý toán học và cơ chế thuật toán cho phép hệ thống tích lũy tri thức mới một cách bền vững mà không làm suy giảm khả năng xử lý trên các phân bố dữ liệu đã ghi nhận trong quá khứ.

\subsection{Mâu thuẫn giữa tính dẻo và tính ổn định}

Mâu thuẫn giữa tính dẻo và tính ổn định phản ánh thách thức nền tảng trong việc thiết kế các hệ thống học thích ứng động. Tính dẻo đại diện cho khả năng của hệ thống trong việc điều chỉnh trạng thái không gian tham số để tiếp thu thông tin mới từ luồng dữ liệu hiện tại \cite{grossberg1987competitive}. Ngược lại, tính ổn định đảm bảo rằng các cấu trúc biểu diễn đã được tối ưu hóa trong quá khứ không bị suy biến. Khi một hệ thống quá nghiêng về tính dẻo, hiện tượng quên lãng thảm khốc sẽ xảy ra, làm mất hoàn toàn khả năng nhận dạng các phân lớp dữ liệu cũ \cite{french1999catastrophic}.

Khung toán học tổng quát của bài toán học liên tục được mô hình hóa thông qua một chuỗi các nhiệm vụ nối tiếp $t \in \{1, 2, \dots, T\}$. Tại mỗi bước thời gian $t$, mô hình tiếp nhận phân bố dữ liệu $\mathcal{D}_t$ và hướng tới việc tối ưu hóa hàm mất mát toàn cục được biểu diễn trong Công thức \ref{eq:continual_learning_general_loss}:

\begin{equation}
\mathcal{L}_{\text{total}}(\theta_t) = \mathbb{E}_{(x, y) \sim \mathcal{D}_t} \left[ \ell(f(x; \theta_t), y) \right] + \Omega(\theta_t, \theta_{1:t-1})
\label{eq:continual_learning_general_loss}
\end{equation}

Trong Công thức \ref{eq:continual_learning_general_loss}, ký hiệu $\mathcal{L}_{\text{total}}(\theta_t)$ đại diện cho tổng giá trị mất mát của mô hình tại thời điểm $t$, $\theta_t$ thể hiện tập hợp tham số trạng thái của mạng tại bước $t$, $\mathbb{E}_{(x,y) \sim \mathcal{D}_t}$ là giá trị kỳ vọng toán học trên phân bố dữ liệu của nhiệm vụ $t$, $\ell(f(x; \theta_t), y)$ thể hiện hàm mất mát đo lường độ sai lệch giữa đầu ra dự báo $f(x; \theta_t)$ và nhãn mục tiêu $y$, và $\Omega(\theta_t, \theta_{1:t-1})$ đại diện cho thành phần phạt số học nhằm bảo toàn tri thức từ các tập tham số lịch sử $\theta_{1:t-1}$.

\subsection{Các hướng tiếp cận lý thuyết chủ đạo trong Học liên tục}

Các phương pháp nghiên cứu học liên tục trong lý thuyết học máy được phân chia thành ba nhóm tiếp cận chính dựa trên cơ chế kiểm soát không gian tham số và bộ nhớ. Nhóm tiếp cận thứ nhất dựa trên ràng buộc tham số nhằm hạn chế sự biến động của các trọng số quan trọng đối với các nhiệm vụ trước đó thông qua thành phần phạt trong hàm mất mát \cite{kirkpatrick2017overcoming}. Nhóm tiếp cận thứ hai dựa trên việc phát lại dữ liệu hoặc tạo mẫu tái tạo nhằm duy trì sự hiện diện của phân bố dữ liệu cũ trong quá trình tối ưu hóa các mẫu dữ liệu mới \cite{lopez2017gradient}. Nhóm tiếp cận thứ ba dựa trên việc mở rộng linh hoạt kiến trúc mạng, trong đó không gian biểu diễn được bổ sung các thành phần mới khi tiếp nhận luồng dữ liệu chứa thông tin chưa từng xuất hiện.

Cấu trúc phân loại các hướng tiếp cận lý thuyết chủ đạo trong bài toán học liên tục được tổng hợp trong Hình \ref{fig:continual_learning_taxonomy}.

\begin{figure}[h!]
\centering

\begin{tikzpicture}[
    every node/.style={
        font=\normalsize,
        align=center
    },
    box/.style={
        rectangle,
        rounded corners=2pt,
        thick,
        fill=white,
        draw=black,
        minimum height=1.35cm,
        inner xsep=6pt,
        inner ysep=6pt
    },
    arrow/.style={
        ->,
        thick,
        gray!60
    }
]

% =========================================================
% Luồng dữ liệu
% =========================================================
\node[
    box,
    text width=2.8cm
] (data) at (1.5,1.8) {
    Luồng dữ liệu\\
    $\mathcal{D}_1,\mathcal{D}_2,\dots,\mathcal{D}_T$
};

% =========================================================
% Cơ chế điều phối
% =========================================================
\node[
    box,
    text width=2.8cm
] (controller) at (5.0,1.8) {
    Cơ chế điều phối\\
    Học liên tục
};

% =========================================================
% Các nhánh
% =========================================================

% Ràng buộc tham số
\node[
    box,
    text width=4.2cm
] (regularization) at (9.8,4.0) {
    Ràng buộc tham số\\
    (Regularization-based)
};

% Phát lại dữ liệu
\node[
    box,
    text width=4.2cm
] (replay) at (9.8,1.8) {
    Phát lại dữ liệu\\
    (Replay-based)
};

% Động hóa kiến trúc
\node[
    box,
    text width=4.2cm
] (architecture) at (9.8,-0.4) {
    Động hóa kiến trúc\\
    (Architecture-based)
};

% =========================================================
% Connections
% =========================================================

\draw[arrow]
    (data.east) -- (controller.west);

\draw[arrow]
    (controller.east) -- (regularization.west);

\draw[arrow]
    (controller.east) -- (replay.west);

\draw[arrow]
    (controller.east) -- (architecture.west);

\end{tikzpicture}

\caption{Cấu trúc phân loại các hướng tiếp cận lý thuyết trong bài toán học liên tục.}
\label{fig:continual_learning_taxonomy}

\end{figure}

Trong Hình \ref{fig:continual_learning_taxonomy}, luồng dữ liệu biến đổi theo thời gian được tiếp nhận và xử lý qua khối điều phối trung tâm nhằm cân bằng tính dẻo và tính ổn định. Ba nhánh tiếp cận lý thuyết tương ứng với ba chiến lược tối ưu hóa không gian tham số, bộ nhớ đệm và kiến trúc mô hình.

Mỗi hướng tiếp cận lý thuyết mang lại các đặc tính tính toán riêng biệt trong môi trường xử lý trực tuyến. Phương pháp ràng buộc tham số có ưu điểm là giữ nguyên kích thước mô hình cố định, nhưng chịu giới hạn về dung lượng biểu diễn khi số lượng nhiệm vụ gia tăng. Phương pháp phát lại dữ liệu giúp duy trì chính xác ranh giới phân định của các phân lớp cũ, nhưng đòi hỏi dung lượng bộ nhớ lưu trữ bổ sung cho các tập dữ liệu quá khứ. Phương pháp thích ứng kiến trúc động cho phép mô hình mở rộng dung lượng biểu diễn một cách tự nhiên theo độ phức tạp của dữ liệu, tuy nhiên đòi hỏi cơ chế kiểm soát sự tăng trưởng để tránh bùng nổ tài nguyên tính toán v.v...

\section*{Kết luận Chương 2}
\addcontentsline{toc}{section}{Kết luận Chương 2}

Chương \ref{chap:c2} đã trình bày tổng quan hệ thống cơ sở lý thuyết về mạng Khí nơ-ron tăng trưởng và các mở rộng toán học nhằm giải quyết bài toán tự động phân loại khuyết tật trong kiểm tra chất lượng màn hình thiết bị di động. Nội dung chương đã phân tích chi tiết nguyên lý vận hành của thuật toán Khí nơ-ron tăng trưởng truyền thống, đồng thời chỉ ra các hạn chế kỹ thuật liên quan đến hiện tượng bùng nổ số lượng nút nơ-ron và nguy cơ sáp nhập sai các cụm khuyết tật riêng biệt do nhiễu không gian. Bên cạnh đó, khung toán học của bài toán học liên tục cũng được hệ thống hóa nhằm làm rõ thách thức về mâu thuẫn giữa tính dẻo và tính ổn định trong môi trường xử lý luồng dữ liệu biến đổi.

Các cải tiến toán học cốt lõi của mô hình EAGNG được tổng hợp nhằm tạo tiền đề khắc phục các tồn tại kỹ thuật trên thông qua bốn cơ chế chính:
(i) mô hình hóa quy luật suy giảm tốc độ học và mức tích lũy sai số theo độ tuổi nút nhằm duy trì tính ổn định của cấu trúc tri thức;
(ii) thiết lập vùng đệm bộ nhớ tích lũy bằng chứng kết hợp phương pháp ước lượng mật độ nhân Gauss để hoãn chèn nút đối với các mẫu dữ liệu ngoại lai;
(iii) đánh giá mật độ không gian cục bộ nhằm kiểm soát sự tăng trưởng số lượng nút tại các vùng dữ liệu có mật độ cao;
(iv) áp dụng lý thuyết hình học bao không gian để kiểm soát điều kiện tạo cạnh nối giữa các cụm nơ-ron.
Bên cạnh đó, việc tích hợp quy tắc phân phối thống kê ba Sigma cho Pha một và thuật toán thành phần liên thông trên đồ thị cho Pha hai tạo nên cơ chế hoàn chỉnh cho phép lọc bỏ ảnh bình thường và phân nhóm chính xác các dạng lỗi bề mặt.

Toàn bộ các nguyên lý toán học và mô hình lý thuyết được tổng hợp trong Chương \ref{chap:c2} đóng vai trò là nền tảng khoa học trực tiếp cho việc thiết kế kiến trúc chi tiết của hệ thống EAGNG. Trên cơ sở lý thuyết đã được thiết lập, Chương \ref{chap:c3} sẽ trình bày chi tiết phương pháp đề xuất, cấu trúc thuật toán tổng quát, quy trình xử lý dữ liệu luồng thực tế và cách thức hiện thực hóa hệ thống phân loại khuyết tật màn hình điện thoại.